\chapter{검증 결과(6): 조정형 Stage C와 R2 통합 회로}

\section{조정형 Stage C 설계(E22)}

R2 기본형에서는 400 V 버스가 V\textsubscript{A}를 그대로 따르므로 Stage C 입력이 270–504 V로 1.87배 바뀐다. 270 V는 설계 입력 하한이며, 180s 팩의 부하를 건 최저 버스(268.77 V; A25R 등가 268.88 V, E46)는 이보다 낮아 그 입력에서 48 V 서비스의 정방향·회생·듀티 한계·전압 품질을 다시 판정해야 한다(7장 표 각주 a). 기존 고정이득 CLLC(권선비 400/48)에 이 전압을 넣으면 48 V 출력이 32.4–60.5 V가 되어 무제한 운전 범위(36–52 V)\cite{b10}를 벗어난다. 또 Stage C는 출력전압과 배터리 충전전류를 스스로 조정하고, 48 V 3 kW와 12 V 1.5 kW의 동시 공급(약 5 kW)과 48 V 회생(−3 kW)을 받아야 한다.

공진형 변환기 하나로 이 이득 범위를 양방향으로 내려면 자화 인덕턴스를 작게 두어 순환전류가 커지므로, 전단 조정기와 공진점 고정 주파수 CLLC 직류변압기(DCX)의 2단 구성을 택했다\cite{b17,b18}. 전단은 E15의 손실 규칙으로 두 후보를 비교했다. A(동기 벅 + CLLC 4.8:1, V\textsubscript{I} 173–250 V, 스위치 위치 10, 인덕터 204 µH·첨두 29.8 A·0.120 J)는 48 V 1.5 kW 평균 손실 47.1 W(재설계 뒤 50.8 W), 5 kW 최악 166.8 W(96.8 \%; 10.2절 재설계 뒤 170.3 W·96.7 \%, E46, 모델 값이며 측정 효율이 아님)이고, B(4스위치 벅부스트, V\textsubscript{I} 300–433 V, 위치 12, 236 µH·첨두 19.1 A·0.057 J)는 49.3 W, 154.8 W(97.0 \%)다. 차가 5 \% 이하(재설계 뒤 3.2 \%)면 모드 경계가 없고 위치가 적은 쪽을 고르는 규칙에 따라 A를 택했다.


선택한 구성은 벅 100 kHz(L\textsubscript{b} 204 µH, C\textsubscript{I} 20 µF), CLLC 권선비 4.8, 공진 2.19 µH·184.9 nF(250 kHz), 출력 커패시터 2.2 mF다. CLLC는 보내는 쪽 브리지만 구동하고 받는 쪽은 동기정류한다(두 브리지를 같은 위상으로 구동하면 중부하에서 영전압 스위칭이 깨진다). 직류변압기의 이득은 −3–5 kW에서 0.975–1.020이었고(출력용량 레그를 넣기 전 직류변압기 모델의 표본 범위), 제어(2.6절)는 평균모델 시험 8종을 모두 통과했다(표 \ref{t:e22c}, 그림 \ref{f:e22}).

{\footnotesize
\begin{longtable}{L{\dimexpr 0.3860\linewidth-2\tabcolsep\relax}L{\dimexpr 0.2632\linewidth-2\tabcolsep\relax}L{\dimexpr 0.3509\linewidth-2\tabcolsep\relax}}
\caption{조정형 Stage C의 제어 시험(E22, 평균모델, 입력 270·330·400·460·504 V)}\label{t:e22c}\\
\toprule
\tabhead \textbf{시험} & \textbf{판정 기준} & \textbf{결과} \\
\midrule
\endfirsthead
\toprule
\tabhead \textbf{시험} & \textbf{판정 기준} & \textbf{결과} \\
\midrule
\endhead
\bottomrule
\endlastfoot
T1 정상 조정(입력 5점 × 전력 4점) & 오차 ≤ 0.1 \% & 0.000 \% \\
T2 계약 부하 변화 4 kW/ms & 48 V ±2 \% & 최대 0.49 \% \\
T3 계단(배터리 없음, 0.1 ↔ 5 kW, −3 ↔ +5 kW) & 36–52 V, 출력 100 µs 평균 ≤ 136.5 A & 40.64–50.91 V, 128.0 A(2.5 ms 뒤 104.2 A) \\
T4 버스 계단 ±50 V/10 µs & 48 V ±2 \% & 최대 0.59 \% \\
T5 배터리 연결(충전 제한 30 A) & 36–52 V, 1 ms 뒤 충전 ≤ 31.5 A & 43.28–48.98 V, 30.0 A \\
T6 방향 전환 −3 → +3 → −3 kW & 48 V ±2 \% & 최대 0.22 \% \\
T7 입력·버스 임피던스(Middlebrook) & |Z\textsubscript{bus}| ≤ |Z\textsubscript{in}|/2 & 최소 여유 49.1 dB \\
T8 경부하 두 브리지 구동(10.2절) & 계단 36–52 V, ±1 kW 램프 ±2 \% & 46.74–48.29 V, 최대 0.21 \% \\
\end{longtable}
}
\srcnote{주: 계단 직후의 순간 전류는 제어로 막을 수 없는 커패시터 사이 전하 이동이라, 전류 제한은 100 µs 이동평균으로 판정하고 순간값(최대 135 A)은 2차 소자 펄스 정격(329 A\cite{b15})과 비교했다.}

\begin{figure}[!htb]
\centering
\includegraphics[width=16.0cm]{\repo/validation/2026-10-25-stage-c-regulated/E22/E22_stage_c_regulated.png}
\caption{조정형 Stage C의 후보 손실, 직류변압기 이득과 48 V 계단 응답(E22)}\label{f:e22}
\srcnote{주: (a) 후보 구조 손실(재설계 전), (b) 직류변압기 이득과 보내는 쪽 영전압 스위칭 여유, (c)–(e) 48 V 계단 응답(입력 270·400·504 V), (f) Stage C 입력 임피던스와 R2 기본형 버스 임피던스.}
\end{figure}

\section{조정형 Stage C의 회로 확인(TB14, TB15)}

소자 수준 회로(Simscape Electrical, 2.5 ns)의 확인은 세 번 걸렸고, 두 번의 실패는 원인을 검증 모델로 재현한 뒤 고쳤다(표 \ref{t:tb14}). \textbf{TB14 1차}는 다섯 사례 모두 켜짐 직전 소자 전압이 51.8–68.5 V로 남아 영전압 스위칭에 실패했다. 1차가 보낼 때는 레그가 데드타임 안에 전이를 끝내지 못했고, 2차가 보낼 때는 전이를 마친 노드가 공진전류 반전으로 되돌아갔다. 처음 스위칭 모델은 데드타임 동안 브리지 전압을 다이오드가 바로 정해 이를 볼 수 없었으므로, 레그 노드와 출력용량을 넣어 회로 결과를 재현하고(5 kW에서 67.9 V 대 회로 66.7–68.5 V) 출력용량 ±20 \%에서도 통과하는 창으로 자화 인덕턴스를 62.2 µH에서 33 µH로, 데드타임을 1차 100 ns·2차 40 ns(창 90–110·25–55 ns의 가운데)로 정했다(1.5 kW 손실 +3.7 W). \textbf{TB14 2차}에서는 받는 쪽 1차 브리지를 다이오드 정류로만 두면 1차 레그 출력용량 때문에 0.63 kW보다 가벼운 회생에서 V\textsubscript{I}가 n·V\textsubscript{o}보다 3 \% 넘게 높아야 했다(ngspice가 8 \% 안에서 재현, 초경부하 LLC의 출력 상승\cite{b19}과 같은 계열). 270 V 버스에서는 벅 듀티가 상한에 붙어 조정이 깨질 수 있어, 0.6 kW 아래에서 받는 쪽 브리지도 같은 위상으로 구동했다(2차 데드타임 90 ns, TB15).

\Needspace{8\baselineskip}
{\footnotesize
\begin{longtable}{L{\dimexpr 0.1491\linewidth-2\tabcolsep\relax}L{\dimexpr 0.2632\linewidth-2\tabcolsep\relax}L{\dimexpr 0.0877\linewidth-2\tabcolsep\relax}L{\dimexpr 0.5000\linewidth-2\tabcolsep\relax}}
\caption{조정형 Stage C 회로 모델의 판정(TB14, TB15)}\label{t:tb14}\\
\toprule
\tabhead \textbf{실행} & \textbf{설계} & \textbf{판정} & \textbf{결과} \\
\midrule
\endfirsthead
\toprule
\tabhead \textbf{실행} & \textbf{설계} & \textbf{판정} & \textbf{결과} \\
\midrule
\endhead
\bottomrule
\endlastfoot
TB14 1차 & 자화 62.2 µH, 데드타임 100 ns & 0/5 & 켜짐 직전 소자 전압 51.8–68.5 V(기준 레일의 5 \%)로 영전압 스위칭 불통과 \\
TB14 2차 & 자화 33 µH, 데드타임 100·40 ns. 역방향 경부하 F·G 추가 & 6/7 & 켜짐 직전 ≤ 0.19 V(A–F). F·G 끝 이득 0.910·0.911(예상 약 0.91, TB14 판정 항목 아님). 270 V 사례 G는 켜짐 직전 2차 40.0 V와 E22 끝 전류 대조로 불통과(벅 듀티 상한, 보내는 쪽 전환 41번) \\
TB15 & TB14 2차 + 경부하 두 브리지 구동(2차 90 ns) & 5/5 & F·G 끝 이득 1.001·1.000, 켜짐 직전 ≤ 0.22 V, 모드 전환 사례 H(400 V, +1→−1 kW)의 48 V 47.79–48.07 V(E22와 최저 차 0.164 V, 끝 차 0.0035 V, 원 기준 통과) \\
\end{longtable}
}
\srcnote{주: 판정 묶음은 48 V 범위, 출력·소자 전류, 켜짐 직전 소자 전압(레일의 5 \% 이하), 끝 직류변압기 이득(1 ± 3 \%), E22 대조다. 출력 정격 105 A(과도 130 A, 2 ms)는 지령·100 µs 평균 기준이다(저장 회로 요약의 100 µs 평균 최대: TB15 106.96 A, TB14 사례 E 129.98 A). 영전압 스위칭은 각 브리지 대표 레그의 지정 정상 시간창에서만 확인했다.}

남은 것은 두 가지다. 데드타임 창의 폭이 20–30 ns라 하드웨어의 비선형 출력용량과 게이트 지연 편차로 다시 정해야 하고, 보내는 쪽 하나로 운전하는 역방향 하한(0.6–0.74 kW)에서 이득이 0.96–0.97로 기준 1 ± 3 \%를 최대 약 1 \%p 벗어난다(48 V 설정의 조정은 유지). 또 \textbf{CV 설정 범위 36–52 V 전체는 보장하지 않는다.} 입력 270 V·출력 52 V·−650 W·두 브리지 출력용량 +20 \%에서, 큰 회생에서 내려온 단측 송신 정상점은 필요 듀티 0.980이 한도 0.97을 넘고(공칭 출력용량 0.969, 268.77 V에서는 공칭도 0.974), 양측 구동이면 0.922로 가능하다(E46; 정상점 계산이며 폐루프 발산이 아니고, 600–700 W 모드 히스테리시스 안에서는 이전 상태가 결과를 바꿈). 48 V 공칭 계약 램프와 대표 회생 사례는 기준을 만족했고, 180s 최저 버스(268.77 V)에서 48 V 설정의 평균모델 램프 네 전이는 최대 편차 0.49 \%(포화 없음), 계단 40.23–50.91 V였다(E46, 부분 확인). 36–52 V 전 설정·저입력·공차·두 모드 이력의 조합은 별도 검증 대상이다(19장).

\section{R2 통합 회로(E23, TB16, TB17)}

지금까지 각 단은 다른 단의 등가로 두고 확인했다(Stage B: Stage A 정적 등가 V\textsubscript{bat}/2 + 13.2 mΩ, Stage C 정전력 부하. 조정형 Stage C: 이상 버스 전원). 세 단을 한 회로로 이은 이유는 세 가지다. R2 기본형에는 버스를 붙잡는 조정기가 없어 Stage C(정전력 5 kW)와 400 V 부하가 Stage A 출력저항·C\textsubscript{A}·감쇠기에 직접 물리고 안정성은 Stage C 입력 임피던스와 버스 임피던스의 관계로 정해진다\cite{b20}, \cite{b21}. 창 조정 모듈의 H-브리지 개방 강등 동안 Stage C 입력은 460 → 380 → 504 V로 최대 약 130 V/ms 움직인다. 고장 직후 Stage A 경로전류는 약 50 A에서 수 A로 떨어지고 잠시 역방향이 되는데, 영전류 스위칭은 정상 운전점(TB06·TB07)에서만 판정했었다.

TB16·TB17은 150 kHz Stage A로 했다. TB16은 배터리 → Stage A 모듈 1 → V\textsubscript{A} 노드(400 µF) → 배선 → 버스(400 µF ∥ 감쇠기) → 정전력 부하 + 조정형 Stage C(TB15) 회로이고, TB17은 TB13 회로의 이상 V\textsubscript{A} 전원과 정전력 부하를 실제 Stage A·Stage C 회로로 바꿨다(표 \ref{t:r2int}, 그림 \ref{f:e23}).

\Needspace{8\baselineskip}
{\footnotesize
\begin{longtable}{L{\dimexpr 0.1930\linewidth-2\tabcolsep\relax}L{\dimexpr 0.2807\linewidth-2\tabcolsep\relax}L{\dimexpr 0.0877\linewidth-2\tabcolsep\relax}L{\dimexpr 0.4386\linewidth-2\tabcolsep\relax}}
\caption{R2 통합 회로의 판정(E23, TB16, TB17)}\label{t:r2int}\\
\toprule
\tabhead \textbf{시험} & \textbf{조건} & \textbf{판정} & \textbf{결과} \\
\midrule
\endfirsthead
\toprule
\tabhead \textbf{시험} & \textbf{조건} & \textbf{판정} & \textbf{결과} \\
\midrule
\endhead
\bottomrule
\endlastfoot
E23 기본형(평균모델) & 540·806·907 V. Stage C 0.1 → 5 kW 계단, 400 V 부하 15 → 25 kW(10 µs), Stage C 5 → −3 kW 계단 & 3/3 & 48 V 44.31–48.18 / 48.00–48.03 / 47.75–50.91 V. 버스 변동 0.43–0.83 V. Middlebrook 여유 최소 49.3 dB \\
TB16 기본형(회로) & 같음 & 3/3 & 48 V 44.16–48.19 / 48.00–48.04 / 47.73–50.73 V(E23과 0.18 V 안). Stage A 턴오프 전류/피크 ≤ 2.27 \%, Stage C 켜짐 직전 ≤ 0.21 V \\
E23 강등(평균모델) & E17 버스 궤적, Stage C 입력 5 kW & 2/2 & 48 V 47.992–48.012 V(±0.03 \%) \\
TB17 강등(회로) & R2-W460 1008·930 V, H-브리지 개방, 400 V 부하 20 kW + Stage C 5 kW & 2/2 & 버스 409.4–503.9 / 380.3–464.9 V 동안 48 V 47.995–48.011 V. DAB 첨두 66.4·64.5 A(≤ 68.7 A) \\
\end{longtable}
}
\srcnote{주: 48 V 범위 기준은 계단 사례 36–52 V, 계약 과도·강등 ±2 \%. TB16은 10 µs 평균(벅 한 주기, Stage A 리플 6주기)으로 판정했다.}

\begin{figure}[!htb]
\centering
\includegraphics[width=15.5cm]{\repo/validation/2026-10-25-r2-integrated/E23/E23_degrade.png}
\caption{R2-W460 H-브리지 개방 강등 중 400 V 버스와 조정형 Stage C의 48 V(E23)}\label{f:e23}
\srcnote{주: \textbf{Stage C 보호 일시 정지(9.6절)를 도입하기 전의 보호 순서}와 150 kHz Stage A로 계산한 결과다. 최종 보호 순서에서는 링크 붕괴·소프트 투입 동안 Stage C가 최대 12.4 ms 멈추고 48 V는 48 V 배터리 전압을 따른다(10.4절). 같은 조건의 회로(TB17)에서 48 V는 평균모델과 0.004 V 안이었다.}
\end{figure}


\textbf{설계에 쓴 등가가 회로와 맞았다.} TB17의 강등은 TB13과 같았고(순시 경로·필터 −3.3 – +0.2 A, DAB 첨두 ≤ 0.1 A 차), 실제 Stage C 입력이 등가보다 최대 2.4 \% 컸지만 상태 전이 시각은 바뀌지 않았다. 유효 범위는 V\textsubscript{A} 노드 공진(150 kHz, C\textsubscript{A} 400 µF에서 주기 약 30 µs)보다 느린 외란이다. 강등이 48 V에 준 영향은 ±0.01 V 수준이었다(150 kHz Stage A, Stage C 일시 정지를 넣기 전의 보호 순서). \textbf{새 한계는 Stage A 경부하 영전류 스위칭이었다.} 1008 V 강등 뒤 Stage C만 남으면 모듈 출력은 약 10 A로, 스위칭 모델(E1)의 턴오프 전류/피크가 5.40–5.47 \%로 기준 5 \%를 넘었고 회로는 4.44 \%로 통과했지만 여유가 작았다. 11장에서 해결했다.

\section{보호가 Stage C에 주는 요구와 최종 설계 반영}

보호 상태기계를 SALS·부품 정격과 연동하면서(9.5·9.6절) 조정형 Stage C에 회생 차단, 보호 일시 정지, 이중 고장 부하 정지의 세 요구가 더해졌다(값은 표 \ref{t:req}, 근거 E29, E38, TB23, TB31). 일시 정지는 H-브리지 고장 강등에서만 하며 M0 등급(감발하지 않는 부하)의 유일한 예외이고, 강등 0.5 ms 뒤 다시 켠다. 제어기 구현과 하드웨어 응답은 19장의 확인 과제다.


\textbf{일시 정지 동안의 48 V(E43).} 최종 보호 순서에서는 H-브리지 고장 강등의 링크 붕괴·소프트 투입 동안 Stage C가 최대 12.4 ms 멈추므로 TB17의 48 V 결과를 그대로 쓸 수 없고, 이 구간은 TB31과 E38 계산에 없었다. 48 V 노드를 출력 커패시터 2.2 mF, 배터리 등가, 정전력 부하 5 kW로 두고 Stage C 출력을 12.4 ms 동안 0으로 둔 보수 계산을 했다. E22의 배터리 등가(OCV 44–47.9 V, 20 mΩ)에서 48 V는 41.6–45.7 V로 ISO 21780 무제한 운전 범위(36–52 V) 안이지만, ±2 \% 조정은 적용되지 않는다. 36 V 이상을 지키려면 OCV 44 V에서 내부저항 ≤ 57.6 mΩ, 방전 전류 ≥ 139 A(12.4 ms)가 필요하다(60 mΩ에서 35.6 V로 불통과). 배터리 인덕턴스·온도·SOC 변화와 Stage C 재개 과도는 넣지 않았으므로, 이것은 검증이 아니라 배터리 요구의 도출이다.

최종 설계값(벅 100 kHz, CLLC 4.8:1, 자화 33 µH, 5 kW)은 6장의 손실 비교에도 넣었다. 대표 부하(48 V 1.5 kW, 버스 410 V)의 조정형 Stage C 손실은 51.2 W로, 같은 규칙의 고정이득 5 kW 직류변압기(28.9 W)보다 약 22 W 크다. 이것이 48 V 단의 조정을 Stage B에서 Stage C로 옮긴 대가다.

A25R로 바꾼 뒤 R2 기본형 버스와 48 V 시험(E23 기본형 3사례)을 다시 계산해 3/3 통과했고 Middlebrook 여유는 최소 49.9 dB(기준 6 dB)였다(E36). A25R과 조정형 Stage C를 함께 넣은 회로는 충전 중 정지 시험(TB28·TB30·TB32·TB33, 12장)에서 확인했다. TB30의 클램프 사건 동안 48 V는 47.978–48.029 V, Stage C 켜짐 직전 소자 전압은 0.21 V 이하였다.
